\documentclass[a4paper]{arrowhead}
\usepackage[T1]{fontenc}
\usepackage{longtable,booktabs,xurl}
\setlength{\emergencystretch}{3em}
\setlength{\parindent}{0pt}
\setlength{\parskip}{5pt}
\begin{document}
\ArrowheadTitle{TestTranslation (Go)}
\ArrowheadType{System Design Description}
\ArrowheadTypeShort{SysDD}
\ArrowheadVersion{5.2.2}
\ArrowheadDate{2026-10-24}
\ArrowheadAuthor{Implementation working draft}
\ArrowheadStatus{DEVELOPMENT}
\ArrowheadContact{Repository maintainers}
\ArrowheadFooter{Eclipse Arrowhead -- native Go support implementations}
\ArrowheadSetup
\section{Design}
\subsection{Structure}
The shared support runtime provides mutual TLS, identity-token binding, ConsumerAuthorization checks, registration, telemetry and the JSON state. \path{testtranslation.go} adds the interface translator, the data model translator and the bridge data plane (route \path{/testtranslation/bridge/{bridgeId}/{path...}}, access \texttt{bridge}).
\subsection{Bridge lifecycle}
\begin{enumerate}
\item \texttt{check-targets}: returns the targets having a \texttt{generic\_https} interface (policy \texttt{CERT\_AUTH} or \texttt{NONE}) with an access address, port and the requested operation.
\item \texttt{initialize-bridge}: stores the bridge (operation method and path template, target interface properties, input and result data model translator descriptors) and answers the bridge interface: \texttt{generic\_https}, policy \texttt{TRANSLATION\_BRIDGE\_TOKEN\_AUTH}, base path \path{/testtranslation/bridge/<bridgeId>}, the operation, and \texttt{dataModels} of the consumer.
\item A bridge call: token and operation check; the body is translated by the input data model translator (\texttt{init-translation}, then \texttt{get-translation-result}); the request (query string included) is forwarded to the target with TestTranslation's identity; a successful result is translated by the result data model translator; \texttt{USED} is reported. An unreachable target reports \texttt{EXTERNAL\_ERROR} and ends the bridge; a failed translation reports \texttt{INTERNAL\_ERROR}.
\item \texttt{abort-bridge/<bridgeId>}: only the system that initialized the bridge may abort it; the bridge stops forwarding (404).
\end{enumerate}
\subsection{Persistence}
Bridges: JSON state (\texttt{bridges}). Translation tasks: memory. No secrets are stored.
\subsection{Important delimitations}
TestTranslation is a test application system for the TranslationManager, not a production translator.
\begin{itemize}
\item One bridge kind: \texttt{generic\_https} to \texttt{generic\_https}, one negotiated operation per bridge; targets with \texttt{CERT\_AUTH} or \texttt{NONE} policies (authorization tokens for token-protected targets are not presented).
\item Data models \texttt{simpleJson} and \texttt{senml} only; translations are synchronous (a task is \texttt{DONE} or \texttt{ERROR} when \texttt{init-translation} returns).
\item Bridges are kept in the system's JSON state; translation tasks only in memory (latest 1000).
\end{itemize}
\end{document}
